\documentclass[10pt]{mypackage}

%\usepackage{mlmodern}
%\usepackage{newpxtext,eulerpx,eucal}
%\renewcommand*{\mathbb}[1]{\varmathbb{#1}}

%\usepackage{homework}
%\usepackage{notes}

\usepackage[ backend=bibtex, style = alphabetic, sorting=ynt ]{biblatex}
\addbibresource{all_references.bib}

\usepackage{parskip}

\fancyhf{}
\fancyhead[R]{Avinash Iyer}
\fancyhead[L]{Analysis on Locally Compact Abelian Groups}
\fancyfoot[C]{\thepage}

\setcounter{secnumdepth}{0}

\begin{document}
\RaggedRight
Throughout, we will write the group law as multiplication, even if we are concerned with $\R$ or $\Z$ where the group law is traditionally written out via addition. Additionally, we will let $L_xf(y) = f\left( x^{-1}y \right)$ denote the left regular representation.
\section{Dual Groups}%
If $G$ is a locally compact abelian group, then all the irreducible representations of $G$ are the characters $\xi\colon G\rightarrow S^{1}$. The set of all characters of $G$ are denoted $\hat{G}$, and we will write (via duality discussed in the section on Pontryagin duality) $\xi(x) = \iprod{x}{\xi}$, where $x\in G$ and $\xi\in \hat{G}$.

Each $\xi\in \hat{G}$ determines a nondegenerate $\ast$-representation of $L_1(G)$ on $\C$ by
\begin{align*}
  \xi(f) &= \int_{}^{} \iprod{x}{\xi}f(x)\:dx.
\end{align*}
Every multiplicative linear functional (also known as a character) on $L_1(G)$ is thus given by integration against a character of the group.

In particular, we may identify $\hat{G}$ with the character space of $L_1(G)$ via the following discussion. If $\Phi\in \left( L_1(G) \right)^{\ast}$ is some character, then it is given by integration against some $\phi\in L_{\infty}\left( G \right)$. For any $f,g\in L_1$, we may then take
\begin{align*}
  \Phi(f) \int_{}^{} \phi(x)g(x)\:dx &= \Phi(f)\Phi(g)\\
                                     &= \Phi\left( f\ast g \right)\\
                                     &= \int_{}^{} \int_{}^{} \phi(y)f\left( yx{-1} \right)g(x)\:dx\:dy\\
                                     &= \int_{}^{} \Phi\left( L_xf \right)g(x)\:dx
\end{align*}
Therefore, $\Phi(f)\phi(x) = \Phi\left( L_xf \right)$ locally almost everywhere. We may define $\phi(x) = \Phi(L_xf)/\Phi(f)$ for every $x$, so that $\phi$ is continuous; by replacing $x$ with $xy$ and $f$ with $L_yf$, we then get
\begin{align*}
  \phi\left( xy \right)\Phi(f) &= \Phi\left( L_{xy}f \right)\\
                               &= \Phi\left( L_xL_yf \right)\\
                               &= \phi(x)\phi\left( L_yf \right)\\
                               &= \phi(x)\phi(y)\Phi(f),
\end{align*}
so that $\phi\left( xy \right)=\phi(x)\phi(y)$. Thus, $\phi\left( x^{n} \right) = \phi\left( x \right)^{n}$, meaning that since $\phi$ is bounded, $\left\vert \phi(x) \right\vert = 1$, or that $\phi$ is a character on $G$.

We have that $\hat{G}$ is an abelian group under pointwise multiplication, has the constant function for $1$ as the identity element, and satisfies
\begin{align*}
  \iprod{x}{\xi^{-1}} &= \iprod{x^{-1}}{\xi}\\
                      &= \overline{ \iprod{x}{\xi} }.
\end{align*}
Endowing $G$ with the topology of uniform convergence on compact subsets. This coincides with the weak* topology on $L_{\infty}(G)$, but since $\hat{G}\cup \set{0}$ is the set of all characters from $L_1(G)$ to $\C$, it follows that $\hat{G}$ is locally compact. Thus, we consider $\hat{G}$ to be a locally compact abelian group in and of itself, and we call it the \textit{dual group} of $G$.
\subsection{Properties and Examples}%

\subsection{Characters of $p$-adic Numbers}%

\section{The Fourier Transform}%
\section{Pontryagin Duality}%
\nocite{folland_abstract_harmonic_analysis}
\printbibliography
\end{document}
